\chapter{Графические интерфейсы на JavaFX}
\label{ch:09}

До сих пор наши программы общались с пользователем через консоль или через
браузер: в главе~\ref{ch:06} мы научились хранить данные в базе, в
главе~\ref{ch:07} подняли приложение на Spring, в главе~\ref{ch:08} разобрали
устройство веб-архитектуры. Но есть большой класс программ, которым браузер не
нужен: терминал банковского операциониста, редактор схем, программа складского
учёта. Они запускаются локально, работают быстро и имеют прямой доступ к файлам
и оборудованию; такие программы называют настольными (\eng{desktop}), и в Java
их пишут на JavaFX. В этой главе вы подключите JavaFX к проекту, разберётесь в
жизненном цикле приложения и графе сцены, научитесь раскладывать элементы по
окну, обрабатывать нажатия, связывать данные с интерфейсом привязками, выносить
разметку в FXML и оформлять её через CSS. Завершает главу сквозной пример~---
«Менеджер студентов», собранный по шаблону MVC.

Если продолжать строительную аналогию, до сих пор мы прокладывали инженерные
коммуникации: трубы, проводку, фундамент. Теперь займёмся тем, что видит и
трогает жилец,~--- дверными ручками, выключателями и обоями. Работа не менее
ответственная: именно по ней пользователь судит обо всей системе.

\section{JavaFX среди графических библиотек Java}

За тридцать лет в Java сменилось три поколения графических библиотек, и
различать их полезно: код всех трёх до сих пор встречается в живых проектах.

\term{AWT} (\eng{Abstract Window Toolkit}, 1995) не рисовал элементы сам: он
просил операционную систему создать «настоящий» системный элемент (такие
компоненты называют \term{тяжеловесными}). Приложение выглядело как родное для
системы, но набор элементов пришлось урезать до общего знаменателя всех
платформ. Это как сеть кофеен, заказывающая стулья у местного столяра в каждом
городе: в Москве сделают один стул, в Токио~--- другой, а если токийский столяр
не умеет делать подлокотники, от них придётся отказаться во всей сети.

\term{Swing} (1998) рисует все элементы сам средствами Java~2D
(\term{легковесные} компоненты): стулья теперь делает собственная фабрика, они
одинаковы везде и могут быть какими угодно. Swing огромен и до сих пор работает
в промышленных системах, но его архитектура родом из девяностых: разметка
задаётся только кодом, стилизация болезненна, анимации почти нет.

\term{JavaFX} (2008, современный вид~--- с 2011~года)~--- третье поколение и
официальная рекомендация для новых настольных приложений: здесь взяли лучшее от
Swing и добавили то, что к тому времени стало нормой в вебе. Различия поколений
сведены в табл.~\ref{tab:09-1}; главные строки~--- две последние, именно они
определяют выбор для нового проекта.

\begin{table}[htbp]
\caption{Три поколения графических библиотек Java}
\label{tab:09-1}
\centering
\begin{tabular}{L{3.6cm}L{3.2cm}L{3.8cm}L{4.2cm}}
\toprule
Характеристика & AWT & Swing & JavaFX \\
\midrule
Год появления & 1995 & 1998 & 2008 \\
Компоненты & тяжеловесные & легковесные & легковесные, ускоренные \\
Внешний вид & как в системе & единый, настраиваемый & единый, настраивается CSS \\
Разметка & только код & только код & код или FXML \\
Привязка данных & нет & нет & свойства и привязки \\
Входит в JDK & да & да & нет, начиная с JDK~11 \\
Статус & устарел & поддерживается & активно развивается \\
\bottomrule
\end{tabular}
\end{table}

Четыре особенности JavaFX стоит назвать отдельно. \term{Граф сцены}
(\eng{scene graph}): интерфейс~--- дерево объектов, а не набор команд рисования;
вы описываете, \emph{что} находится на экране, а не \emph{как} это нарисовать.
\term{Разделение разметки и логики}: интерфейс описывается в XML-файле FXML,
поведение~--- в классе-контроллере. \term{CSS-оформление}: цвета, шрифты и
отступы меняются без перекомпиляции. \term{Свойства и привязки}: значение поля
ввода можно связать с текстом метки, и она обновится сама. Добавьте сюда
аппаратное ускорение через конвейер \eng{Prism} и богатый набор компонентов~---
от кнопки до таблиц, диаграмм и встроенного браузера \texttt{WebView}.

\begin{vrezka}
\textbf{Почему JavaFX больше не входит в JDK.} До Java~10 включительно JavaFX
поставлялся вместе с Oracle JDK. Начиная с JDK~11 (2018) его вынесли в отдельный
проект с открытым исходным кодом~--- OpenJFX. Причин три: JavaFX развивается
быстрее самой Java, и привязка к релизам JDK его тормозила; серверному
приложению графика не нужна и только раздувает дистрибутив; отдельный модуль
легче версионировать. Следствие жёсткое: если просто поставить JDK~21 и написать
\texttt{import javafx.application.Application;}, проект не скомпилируется.
Поэтому следующий раздел пропускать нельзя~--- без него не запустится ни один
пример главы.
\end{vrezka}

\section{Подключение JavaFX через Maven}

Раньше набор отвёрток лежал в общем ящике, который выдавали каждому новому
сотруднику. Теперь ящик легче, а отвёртки заказывают отдельной коробкой и
указывают, где она лежит. Неудобно ровно один раз~--- при настройке проекта.

JavaFX разбит на модули, и подключать все семь не нужно (табл.~\ref{tab:09-2}):
для большинства задач хватает двух верхних строк, поскольку
\texttt{javafx-controls} сам подтягивает \texttt{javafx-graphics} и
\texttt{javafx-base}.

\begin{table}[htbp]
\caption{Модули JavaFX и их назначение}
\label{tab:09-2}
\centering
\begin{tabular}{L{4.2cm}L{10.8cm}}
\toprule
Модуль & За что отвечает \\
\midrule
\texttt{javafx-controls} & элементы управления: \texttt{Button}, \texttt{Label}, \texttt{TableView} \\
\texttt{javafx-fxml} & загрузка интерфейса из файлов разметки FXML \\
\texttt{javafx-graphics} & граф сцены, \texttt{Stage}, \texttt{Scene}, геометрия, анимация \\
\texttt{javafx-base} & свойства, привязки, наблюдаемые коллекции \\
\texttt{javafx-web} & компонент \texttt{WebView}~--- браузер на движке WebKit \\
\texttt{javafx-media} & воспроизведение аудио и видео \\
\bottomrule
\end{tabular}
\end{table}

Сборочный файл проекта показан в листинге~\ref{lst:09-1}. Обратите внимание на
две детали: версия JavaFX вынесена в свойство и меняется одной строкой, а плагин
\texttt{javafx-maven-plugin} нужен затем, чтобы запускать приложение с правильным
путём к модулям~--- вручную его пришлось бы прописывать каждый раз.

\begin{listing}[H]
\begin{minted}{xml}
<?xml version="1.0" encoding="UTF-8"?>
<project xmlns="http://maven.apache.org/POM/4.0.0">
  <modelVersion>4.0.0</modelVersion>
  <groupId>com.example</groupId>
  <artifactId>student-manager</artifactId>
  <version>1.0-SNAPSHOT</version>

  <properties>
    <maven.compiler.release>21</maven.compiler.release>
    <project.build.sourceEncoding>UTF-8</project.build.sourceEncoding>
    <javafx.version>21.0.2</javafx.version>
  </properties>

  <dependencies>
    <dependency>  <!-- тянет graphics и base -->
      <groupId>org.openjfx</groupId>
      <artifactId>javafx-controls</artifactId>
      <version>${javafx.version}</version>
    </dependency>
    <dependency>
      <groupId>org.openjfx</groupId>
      <artifactId>javafx-fxml</artifactId>
      <version>${javafx.version}</version>
    </dependency>
  </dependencies>

  <build><plugins>
    <plugin>  <!-- запуск с правильным module-path -->
      <groupId>org.openjfx</groupId>
      <artifactId>javafx-maven-plugin</artifactId>
      <version>0.0.8</version>
      <configuration>
        <mainClass>com.example.manager.MainApp</mainClass>
      </configuration>
    </plugin>
  </plugins></build>
</project>
\end{minted}
\caption{Сборочный файл pom.xml проекта с JavaFX}
\label{lst:09-1}
\end{listing}

Важная деталь: у зависимостей \texttt{org.openjfx} не указан классификатор
операционной системы. Их собственный сборочный файл сам определяет, на чём вы
работаете, и подставляет нужную сборку~--- поэтому один и тот же
\texttt{pom.xml} годится и для Windows, и для Linux, и для macOS.

Раскладка файлов проекта показана на рис.~\ref{fig:09-1}. Главное соглашение
здесь такое: FXML и CSS кладут в \texttt{src/main/resources}, повторяя структуру
пакетов. Тогда файл, лежащий «рядом» с классом, находится вызовом
\texttt{MainApp.class.getResource("students-view.fxml")}~--- без абсолютных путей
и без догадок о том, куда сборщик положил ресурсы.

\begin{figure}[htbp]\centering
\begin{forest} hierarchy
[student-manager/
  [pom.xml]
  [src/main/
    [java/com/example/manager/
      [MainApp.java]
      [Launcher.java]
      [controller/StudentController.java]
      [model/Student.java]
      [service/StudentService.java]
    ]
    [resources/com/example/manager/
      [students-view.fxml]
      [styles.css]
    ]
  ]
]
\end{forest}
\caption{Структура проекта настольного приложения}\label{fig:09-1}
\end{figure}

Собирается и запускается проект двумя командами, одинаковыми на всех системах:
\texttt{mvn clean compile} скачивает зависимости и компилирует код, а
\texttt{mvn javafx:run} запускает приложение, попутно копируя разметку и стили в
каталог \texttt{target/classes}.

Самая частая ошибка запуска~--- сообщение
\texttt{JavaFX runtime components are missing}: класс с методом \texttt{main}
наследует \texttt{Application}, а модули JavaFX попали в \texttt{classpath}, а не
в путь к модулям (типично при запуске из среды разработки). Надёжное и
совершенно легальное решение~--- отдельный класс-запускалка, который
\texttt{Application} не наследует (листинг~\ref{lst:09-2}): проверка внутри
платформы тогда просто не срабатывает.

\begin{listing}[H]
\begin{minted}{java}
package com.example.manager;

/**
 * Точка входа-обёртка. Не наследует Application, поэтому
 * проверка не срабатывает и приложение стартует даже из IDE.
 */
public class Launcher {
    public static void main(String[] args) {
        MainApp.main(args);
    }
}
\end{minted}
\caption{Класс-запускалка в обход проверки JavaFX}
\label{lst:09-2}
\end{listing}

Вторая по частоте~--- \texttt{Location is required} из загрузчика разметки:
значит, \texttt{getResource()} вернул \texttt{null}, потому что файл лежит в
\texttt{src/main/java} вместо \texttt{src/main/resources}, перепутан путь или
проект не пересобран.

\section{Класс Application и жизненный цикл приложения}

Любое приложение JavaFX начинается с класса, наследующего абстрактный класс
\texttt{javafx.application.Application} (листинг~\ref{lst:09-3}).

\begin{listing}[H]
\begin{minted}{java}
package com.example.manager;

// импорты: Application, Scene, Label, StackPane, Stage

public class HelloApp extends Application {

    // Единственный абстрактный метод класса Application
    @Override
    public void start(Stage primaryStage) {
        Label label = new Label("Здравствуй, JavaFX!");
        StackPane root = new StackPane(label);   // корень графа сцены
        Scene scene = new Scene(root, 400, 200); // сцена 400 на 200

        primaryStage.setTitle("Моё первое окно");
        primaryStage.setScene(scene);
        primaryStage.show();      // без этой строки окна не будет
    }

    public static void main(String[] args) {
        launch(args);             // запуск среды выполнения JavaFX
    }
}
\end{minted}
\caption{Минимальное приложение JavaFX}
\label{lst:09-3}
\end{listing}

Метод \texttt{start} получает готовое окно и наполняет его содержимым: метка
кладётся в панель, панель становится корнем сцены, сцена ставится в окно, окно
показывается. Метод \texttt{main} вызывает унаследованный \texttt{launch}, а тот
запускает целую машину (рис.~\ref{fig:09-2}): порядок шагов на схеме~--- это
порядок, в котором платформа передаёт управление вашим методам.

\begin{figure}[htbp]\centering
\begin{tikzpicture}[font=\small,node distance=5mm,
  st/.style={draw,rounded corners=2pt,align=center,text width=92mm,
    inner sep=4pt,minimum height=9mm}]
\node[st] (a) {\texttt{main()} вызывает \texttt{launch(args)}};
\node[st,below=of a] (b) {Запускается среда выполнения JavaFX; объект
  приложения создаётся через рефлексию};
\node[st,below=of b] (c) {Вызывается \texttt{init()}\\
  \scriptsize поток JavaFX-Launcher};
\node[st,below=of c] (d) {Создаётся первичный \texttt{Stage}, вызывается
  \texttt{start(stage)}\\\scriptsize поток JavaFX Application Thread};
\node[st,below=of d] (e) {Приложение живёт и обрабатывает события до закрытия
  последнего окна или вызова \texttt{Platform.exit()}};
\node[st,below=of e] (f) {Вызывается \texttt{stop()}, среда завершается,
  управление возвращается в \texttt{main()}};
\draw[flowarrow] (a) -- (b);
\draw[flowarrow] (b) -- (c);
\draw[flowarrow] (c) -- (d);
\draw[flowarrow] (d) -- (e);
\draw[flowarrow] (e) -- (f);
\end{tikzpicture}
\caption{Жизненный цикл приложения JavaFX}\label{fig:09-2}
\end{figure}

Три метода жизненного цикла вы можете переопределить, и у каждого своя зона
ответственности (табл.~\ref{tab:09-3}). Обязателен только \texttt{start}: он
объявлен абстрактным.

\begin{table}[htbp]
\caption{Методы жизненного цикла приложения}
\label{tab:09-3}
\centering
\begin{tabular}{L{2.4cm}L{3.4cm}L{3.6cm}L{5.2cm}}
\toprule
Метод & Когда вызывается & В каком потоке & Что в нём делают \\
\midrule
\texttt{init()} & один раз, до создания окна & JavaFX-Launcher & читают настройки, подключаются к базе; создавать окно и сцену нельзя \\
\texttt{start()} & один раз, после \texttt{init()} & JavaFX Application Thread & строят граф сцены и показывают окно \\
\texttt{stop()} & при завершении приложения & JavaFX Application Thread & сохраняют настройки, закрывают соединения \\
\bottomrule
\end{tabular}
\end{table}

Аналогия проясняет запрет из первой строки таблицы. \texttt{init()}~--- это
подготовка на кухне до открытия кафе: продукты заказали, печь включили, зал ещё
закрыт. \texttt{start()}~--- открыли двери и обслуживаем гостей.
\texttt{stop()}~--- закрылись, убрали столы, выключили свет. Готовить салаты в
зале при гостях, то есть создавать окно в \texttt{init()}, нельзя: зала ещё
физически нет.

Весь граф сцены живёт в одном потоке~--- \term{JavaFX Application Thread}. Любое
изменение интерфейса должно происходить в нём; попытка сделать это из другого
потока даёт \texttt{IllegalStateException: Not on FX application thread}. Тот же
принцип действует в Swing и в браузере: иначе пришлось бы синхронизировать
каждое поле каждого узла, и интерфейс стал бы невыносимо медленным. Обратная
сторона правила~--- длинную работу в этом потоке выполнять нельзя, иначе окно
«зависнет»: оно просто не сможет перерисоваться. Как мы разбирали в
главе~\ref{ch:05}, тяжёлые операции выносят в отдельные потоки, а результат
возвращают в интерфейс через \texttt{Platform.runLater()}
(листинг~\ref{lst:09-4}).

\begin{listing}[H]
\begin{minted}{java}
Thread worker = new Thread(() -> {
    String result = downloadFromServer();   // долгая операция

    // Только внутри runLater можно трогать узлы интерфейса
    Platform.runLater(() ->
            statusLabel.setText("Загружено: " + result));
});
worker.setDaemon(true);   // не мешает приложению завершиться
worker.start();
\end{minted}
\caption{Возврат результата фоновой работы в интерфейс}
\label{lst:09-4}
\end{listing}

\begin{oshibka}
\textbf{Частая ошибка: выход через System.exit.} Вызов \texttt{System.exit(0)}
обрывает работу виртуальной машины немедленно и метод \texttt{stop()} не
вызывает. Если в \texttt{stop()} вы сохраняете настройки или закрываете
соединение с базой~--- всё это потеряется. Для завершения приложения JavaFX
пользуйтесь \texttt{Platform.exit()}.
\end{oshibka}

\section{Stage, Scene и граф сцены}

Терминология JavaFX взята из театра, и она объясняет структуру приложения лучше
любой документации. \texttt{Stage}~--- сцена театра, то есть окно операционной
системы: рама, заголовок, кнопки свернуть и закрыть, и здание сцены одно и то же
весь вечер. \texttt{Scene}~--- сценическая картина: содержимое окна, которое
между актами меняют целиком. \texttt{Node}~--- актёры и реквизит.

Первичное окно платформа передаёт вам в \texttt{start()}; остальные вы создаёте
сами оператором \texttt{new}. Настраивают окно методами \texttt{setTitle},
\texttt{setScene}, \texttt{setMinWidth}, \texttt{setResizable}. Отдельного
упоминания заслуживает \texttt{setOnCloseRequest}: его обработчик получает
событие закрытия и может отменить его вызовом \texttt{event.consume()}~--- так
спрашивают у пользователя, действительно ли он хочет выйти с несохранёнными
данными. Метод \texttt{show()} возвращает управление немедленно, а
\texttt{showAndWait()} останавливает выполнение до закрытия окна~--- именно он
нужен диалогам, от которых ждут ответа. \term{Модальность} окна задаёт метод
\texttt{initModality}: \texttt{APPLICATION\_MODAL} блокирует все окна
приложения, \texttt{WINDOW\_MODAL}~--- только окно-владельца.

У сцены обязательно есть \term{корневой узел} (\eng{root}), и он должен быть
типа \texttt{Parent}~--- то есть узлом, способным иметь потомков: обычный
\texttt{Node} корнем быть не может, дерево некуда ветвить. Одна и та же сцена не
может стоять сразу в двух окнах, зато у одного окна сцену можно менять сколько
угодно, и это самый простой способ сделать переход между экранами: строка
\texttt{stage.setScene(mainScene)} мгновенно заменяет экран входа на главное
окно.

\subsection{Граф сцены и иерархия узлов}

\term{Граф сцены} (\eng{scene graph})~--- дерево всех визуальных объектов
приложения; корень дерева~--- тот узел, который вы передали в \texttt{Scene}.
Аналогия~--- содержимое шкафа: шкаф (корневая панель) содержит полки (вложенные
панели), полки~--- коробки, коробки~--- вещи. Чтобы вынести всю верхнюю полку,
достаточно вынести один объект вместе со всем содержимым~--- именно так работает
удаление узла в JavaFX.

На рис.~\ref{fig:09-3} показан граф сцены окна с шапкой, таблицей и панелью
ввода~--- ровно того окна, которое мы соберём в конце главы. Подписи
\texttt{top}, \texttt{center} и \texttt{bottom} обозначают зоны панели
\texttt{BorderPane}, о которых пойдёт речь в следующем разделе.

\begin{figure}[htbp]\centering
\begin{forest} hierarchy
[Scene
  [BorderPane -- корень
    [top: Label «Менеджер студентов»]
    [center: TableView
      [TableColumn «ФИО»]
      [TableColumn «Группа»]
      [TableColumn «Средний балл»]
    ]
    [bottom: VBox
      [GridPane
        [Label «ФИО:»]
        [TextField]
      ]
      [HBox
        [Button «Добавить»]
        [Button «Удалить»]
      ]
      [Label «Студентов в списке: 3»]
    ]
  ]
]
\end{forest}
\caption{Граф сцены окна менеджера студентов}\label{fig:09-3}
\end{figure}

Всё, что можно поместить на сцену, наследует абстрактный класс
\texttt{javafx.scene.Node}. Часть этой иерархии приведена на
рис.~\ref{fig:09-4}; запоминать её целиком не нужно, но четыре уровня из неё
знать обязательно.

\begin{figure}[htbp]\centering
\begin{forest} hierarchy
[Node
  [Canvas]
  [ImageView]
  [Shape
    [Rectangle]
    [Circle]
    [Text]
  ]
  [Parent
    [Group]
    [Region
      [Pane
        [VBox]
        [HBox]
        [StackPane]
        [BorderPane]
        [GridPane]
      ]
      [Control
        [Labeled [Label]]
        [ButtonBase [Button]]
        [TextInputControl [TextField] [TextArea]]
        [ComboBox]
        [ListView]
        [TableView]
      ]
    ]
  ]
]
\end{forest}
\caption{Иерархия узлов графа сцены}\label{fig:09-4}
\end{figure}

\texttt{Node}~--- базовый класс любого визуального объекта; узел, который не
является \texttt{Parent} (кнопка, картинка, круг),~--- лист дерева.
\texttt{Parent}~--- ветвь: у него есть список потомков \texttt{getChildren()}, и
именно \texttt{Parent} передаётся в конструктор сцены. \texttt{Region}~---
\texttt{Parent} с прямоугольной областью, фоном, рамкой и отступами; отсюда
растут все панели компоновки и все элементы управления. \texttt{Control}~---
\texttt{Region} со сменной «шкурой» (\eng{skin}) и поддержкой CSS по умолчанию.

\begin{oshibka}
\textbf{Ловушка: Text и Label~--- разные вещи.} Класс
\texttt{javafx.scene.text.Text}~--- это фигура (\texttt{Shape}), она нужна при
рисовании графики. Для подписи в интерфейсе берут
\texttt{javafx.scene.control.Label}: он наследует \texttt{Control}, умеет
переносить текст по словам, показывать картинку рядом с текстом и связываться с
полем ввода.
\end{oshibka}

С деревом узлов работают обычными методами списка. Вызов
\texttt{getChildren().add(node)} добавляет потомка, \texttt{remove(node)}
удаляет узел вместе со всем поддеревом, а \texttt{getParent()} ведёт вверх. Три правила при этом нарушают чаще всего. Первое: у узла может быть
только один родитель~--- добавили кнопку в \texttt{VBox}, потом её же в
\texttt{HBox}, и она молча исчезнет из \texttt{VBox}; нужны две одинаковые
кнопки~--- создайте два объекта. Второе: менять граф сцены можно только из
потока интерфейса. Третье: порядок в \texttt{getChildren()} определяет порядок
отрисовки~--- свойства \texttt{z-index}, как в вебе, здесь нет, чем позже узел
добавлен, тем он выше. Это стопка прозрачных плёнок у мультипликатора: положили
новую сверху~--- она перекрыла всё, что лежит ниже; менять порядок на лету
помогают \texttt{toFront()} и \texttt{toBack()}.

Отдельная пара, которую регулярно путают,~--- \texttt{setVisible} и
\texttt{setManaged}. Вызов \texttt{setVisible(false)} прячет узел, но панель
компоновки продолжает резервировать под него место, и получается «дырка». Чтобы
узел исчез совсем, нужен ещё \texttt{setManaged(false)}.

\section{Панели компоновки}

Расставить элементы по абсолютным координатам можно, но окно, которое
разваливается при изменении размера или при другом системном шрифте,~--- плохое
окно. \term{Панели компоновки} (\eng{layout panes}) сами вычисляют положение и
размер потомков по заданному правилу.

Аналогия с переездом: у вас есть коробки разных типов. В одну вещи складывают
строго в столбик, в другую~--- в ряд, третья разделена на ячейки-сетку,
четвёртая держит крупные предметы по краям, а лёгкие в центре. Вы не измеряете
каждую вещь линейкой~--- вы выбираете подходящую коробку. Какая коробка для чего
предназначена, показывает табл.~\ref{tab:09-4}.

\begin{table}[htbp]
\caption{Панели компоновки JavaFX}
\label{tab:09-4}
\centering
\begin{tabular}{L{2.6cm}L{6.4cm}L{6.0cm}}
\toprule
Панель & Правило раскладки & Типичное применение \\
\midrule
\texttt{VBox} & потомки в столбик, сверху вниз & форма, боковое меню \\
\texttt{HBox} & потомки в строку, слева направо & панель кнопок \\
\texttt{BorderPane} & пять зон: верх, низ, лево, право, центр & каркас главного окна \\
\texttt{GridPane} & сетка из строк и столбцов & формы «подпись~--- поле» \\
\texttt{StackPane} & все потомки друг на друге, по центру & заглушка «Загрузка» \\
\texttt{FlowPane} & в строку с переносом по ширине & плитка тегов, галерея \\
\texttt{AnchorPane} & привязка к краям контейнера & точная подгонка \\
\bottomrule
\end{tabular}
\end{table}

Начнём с двух самых частых панелей~--- строки и столбца
(листинг~\ref{lst:09-5}). Ключевая строка здесь третья: она объясняет, кто
забирает свободное место при растягивании окна.

\begin{listing}[H]
\begin{minted}{java}
// Строка: поле растягивается, кнопка сохраняет свой размер
HBox searchBar = new HBox(8, searchField, new Button("Найти"));
HBox.setHgrow(searchField, Priority.ALWAYS);
searchBar.setAlignment(Pos.CENTER_LEFT);

VBox root = new VBox(12);            // 12 пикселей между потомками
root.setPadding(new Insets(16));     // отступ от краёв панели
root.getChildren().addAll(new Label("Поиск"), searchBar);

VBox.setMargin(searchBar, new Insets(0, 0, 10, 0));
\end{minted}
\caption{Строка и столбец: HBox и VBox}
\label{lst:09-5}
\end{listing}

Здесь встретились два понятия, которые регулярно путают.
\term{Внутренний отступ} (\eng{padding})~--- расстояние от краёв панели до её
потомков, задаётся самой панелью. \term{Внешний отступ} (\eng{margin})~---
свободное место вокруг одного потомка, задаётся статическим методом панели:
\texttt{VBox.setMargin(child, new Insets(10))}. Класс \texttt{Insets} принимает
либо одно число, либо четыре~--- сверху, справа, снизу, слева.

Панель \texttt{BorderPane} делит окно на пять зон и работает как газетная
полоса: шапка сверху, подвал снизу, колонки по бокам, главный материал в центре.
Центр забирает всё оставшееся место~--- то, что нужно таблице или редактору,~---
а незанятые зоны места не занимают вовсе. Для форм идеальна панель
\texttt{GridPane}: подписи в первом столбце, поля во втором
(листинг~\ref{lst:09-6}). Порядок аргументов метода \texttt{add} путают чаще
всего, поэтому запомните его отдельно: сначала номер столбца, потом номер
строки, и нумерация идёт с нуля.

\begin{listing}[H]
\begin{minted}{java}
GridPane grid = new GridPane();
grid.setHgap(10);                // расстояние между столбцами
grid.setVgap(10);                // расстояние между строками
grid.setPadding(new Insets(20));

// add(узел, номерСтолбца, номерСтроки) - нумерация с нуля
grid.add(new Label("Логин:"), 0, 0);
grid.add(new TextField(), 1, 0);
grid.add(new Label("Пароль:"), 0, 1);
grid.add(new PasswordField(), 1, 1);

// add(узел, столбец, строка, ширина, высота) - в ячейках
grid.add(new Button("Войти"), 0, 2, 2, 1);

// Первый столбец фиксированный, второй растягивается
ColumnConstraints labels = new ColumnConstraints(90);
ColumnConstraints fields = new ColumnConstraints();
fields.setHgrow(Priority.ALWAYS);
grid.getColumnConstraints().addAll(labels, fields);
\end{minted}
\caption{Форма на панели GridPane}
\label{lst:09-6}
\end{listing}

Полезный приём отладки: метод \texttt{setGridLinesVisible(true)} рисует линии
сетки прямо в окне, и сразу видно, где узел «уехал». Панели свободно вкладываются друг в друга: почти любое
реальное окно~--- это \texttt{BorderPane}, внутри которого \texttt{VBox}, внутри
которого \texttt{GridPane} и \texttt{HBox}.

\section{Элементы управления}

Элементы управления~--- панель приборов вашего приложения: спидометр показывает
(\texttt{Label}), руль принимает ввод (\texttt{TextField}), кнопка стартера
запускает действие (\texttt{Button}). Водителю не нужно знать, как устроен
датчик,~--- ему нужно, чтобы стрелка была на понятном месте.

\texttt{Label} показывает текст, не принимает ввод и по умолчанию не получает
фокус; самое сильное его применение~--- вывод результата через привязку, о
которой речь пойдёт ниже. Поля ввода образуют семейство с общим предком
\texttt{TextInputControl}: \texttt{TextField}~--- одна строка,
\texttt{PasswordField}~--- она же с замаскированными символами,
\texttt{TextArea}~--- много строк; общие методы у них одни и те же:
\texttt{getText()}, \texttt{setText()}, \texttt{clear()}. Типичные приёмы работы
собраны в листинге~\ref{lst:09-7}: обратите внимание на разницу между реакцией
на клавишу Enter и реакцией на каждое изменение текста.

\begin{listing}[H]
\begin{minted}{java}
TextField field = new TextField();
field.setPromptText("Иванов");   // серая подсказка в пустом поле
field.setPrefColumnCount(20);    // примерная ширина в символах

// Реакция на Enter внутри поля
field.setOnAction(event ->
        System.out.println("Введено: " + field.getText()));

// Реакция на каждое изменение текста
field.textProperty().addListener((obs, oldText, newText) ->
        System.out.println(oldText + " -> " + newText));

// Поле, в которое можно ввести только цифры
field.setTextFormatter(new TextFormatter<>(change ->
        change.getControlNewText().matches("\\d*") ? change : null));
\end{minted}
\caption{Поля ввода и реакция на изменение текста}
\label{lst:09-7}
\end{listing}

\begin{vrezka}
\textbf{Важно.} \texttt{PasswordField} лишь маскирует символы при отрисовке:
\texttt{getText()} возвращает настоящий введённый текст, а не строку из
точек,~--- никакой дополнительной защиты этот класс не даёт.
\end{vrezka}

Кнопка \texttt{Button} наследует цепочку
\texttt{ButtonBase}~$\to$~\texttt{Labeled}~$\to$~\texttt{Control}, поэтому умеет
всё, что умеет подпись, и вдобавок реагирует на нажатие. Главный её метод~---
\texttt{setOnAction}; \texttt{setDefaultButton(true)} заставляет кнопку
срабатывать по клавише Enter в любом месте окна, \texttt{setCancelButton(true)}
--- по клавише Escape, \texttt{setDisable(true)} отключает кнопку. Родственники
кнопки различаются тем, сколько вариантов можно выбрать одновременно:
\texttt{CheckBox}~--- независимый флажок; переключатели \texttt{RadioButton}
объединяют в \texttt{ToggleGroup}, и тогда выбранным остаётся ровно один;
\texttt{ComboBox} и \texttt{ListView} показывают список значений и экономят
место на экране.

\subsection{Таблица данных TableView}

\texttt{TableView}~--- самый мощный и самый востребованный компонент JavaFX. Он
состоит из трёх частей: самой таблицы \texttt{TableView<T>}, типизированной
классом строки; столбцов \texttt{TableColumn<T, V>}, где \texttt{T}~--- тип
строки, а \texttt{V}~--- тип значения в ячейке; и источника данных
\texttt{ObservableList<T>}. Ключевой вопрос: откуда столбец берёт значение? За
это отвечает \term{фабрика значений} (\eng{cellValueFactory}), и задать её можно
двумя способами (листинг~\ref{lst:09-8}).

\begin{listing}[H]
\begin{minted}{java}
TableView<Student> table = new TableView<>();

// Способ 1: по имени свойства, через рефлексию.
// Ищет метод fullNameProperty(), а если нет - getFullName()
TableColumn<Student, String> nameColumn = new TableColumn<>("ФИО");
nameColumn.setCellValueFactory(new PropertyValueFactory<>("fullName"));

// Способ 2: лямбда - типобезопасно, ошибку поймает компилятор
TableColumn<Student, String> groupColumn = new TableColumn<>("Группа");
groupColumn.setCellValueFactory(d -> d.getValue().groupProperty());

table.getColumns().addAll(nameColumn, groupColumn);
table.setItems(service.getStudents());       // связь со списком
table.setPlaceholder(new Label("Список пуст"));

// Реакция на выбор строки
table.getSelectionModel().selectedItemProperty()
        .addListener((obs, was, now) -> {
            if (now != null) System.out.println(now.getFullName());
        });
\end{minted}
\caption{Таблица, столбцы и источник данных}
\label{lst:09-8}
\end{listing}

Про \texttt{PropertyValueFactory} нужно знать главное: он работает через
рефлексию и ищет метод по соглашению об именах, поэтому опечатка в строке
\texttt{"fullName"} не вызовет ошибки компиляции~--- столбец просто останется
пустым. Поэтому в новом коде предпочтительнее лямбда; знать оба способа
обязательно, так как первый встречается повсеместно в существующих проектах.

Если фабрика значений отвечает на вопрос «откуда взять значение», то
\term{фабрика ячеек} (\eng{cellFactory}) отвечает на вопрос «как его показать».
В листинге~\ref{lst:09-9} она округляет средний балл до двух знаков после
запятой; заодно здесь видно, зачем нужен вызов \texttt{asObject()}~--- числовое
свойство работает с примитивом, а таблице нужен объект-обёртка.

\begin{listing}[H]
\begin{minted}{java}
TableColumn<Student, Double> gradeColumn =
        new TableColumn<>("Средний балл");

// ОТКУДА берётся значение
gradeColumn.setCellValueFactory(
        d -> d.getValue().averageGradeProperty().asObject());

// КАК оно выглядит
gradeColumn.setCellFactory(column -> new TableCell<Student, Double>() {
    @Override
    protected void updateItem(Double value, boolean empty) {
        super.updateItem(value, empty);
        setText(empty || value == null
                ? null
                : String.format("%.2f", value));
    }
});
\end{minted}
\caption{Форматирование значения в ячейке}
\label{lst:09-9}
\end{listing}

Писать собственное окно для типового сообщения не нужно~--- есть готовый класс
\texttt{Alert} с типами \texttt{INFORMATION}, \texttt{WARNING}, \texttt{ERROR},
\texttt{CONFIRMATION} и \texttt{NONE}. Текст задают методами
\texttt{setHeaderText} и \texttt{setContentText}, а показывают диалог вызовом
\texttt{showAndWait()}, который ждёт ответа и возвращает
\texttt{Optional<ButtonType>}; пример разбора этого ответа есть в
листинге~\ref{lst:09-16}. Наконец, пакет \texttt{javafx.scene.chart} строит
диаграммы: гистограмме \texttt{BarChart} нужны пара осей и ряд
\texttt{XYChart.Series} с данными. Но в отличие от таблицы диаграмма на список
не подписана: после изменения данных её наполняют заново.

\section{События, свойства и привязка данных}

Событие~--- объект, описывающий, «что произошло». Все они наследуют
\texttt{javafx.event.Event}, у которого есть источник (\texttt{getSource()}),
цель (\texttt{getTarget()}), тип (\texttt{getEventType()}) и метод
\texttt{consume()}~--- «событие обработано, дальше не передавать». Самые частые
классы событий: \texttt{ActionEvent} (нажата кнопка, выбран пункт меню, нажат
Enter в поле), \texttt{MouseEvent}, \texttt{KeyEvent} и \texttt{WindowEvent}.

Обработчик события~--- реализация функционального интерфейса
\texttt{EventHandler} с единственным методом \texttt{handle}, поэтому его пишут
лямбда-выражением или ссылкой на метод (глава~\ref{ch:06}). Помимо «основного
действия» \texttt{setOnAction} есть и низкоуровневые обработчики:
\texttt{setOnMouseClicked} даёт координаты и число щелчков, а
\texttt{setOnKeyPressed}~--- код клавиши и состояние модификаторов.

Когда пользователь щёлкает по кнопке внутри панели внутри окна, событие не
появляется «прямо на кнопке». Оно проходит путь по дереву графа сцены, и
проходит его дважды (рис.~\ref{fig:09-5}). Аналогия: клиент подал жалобу на
конкретного сотрудника. Сначала бумага спускается сверху вниз~--- от директора
через начальника отдела к сотруднику; это \term{фаза захвата}
(\eng{capturing}). Затем, если вопрос не решён на месте, поднимается обратно;
это \term{фаза всплытия} (\eng{bubbling}). Любой на этом пути может написать
«вопрос закрыт» и остановить движение бумаги~--- это и есть \texttt{consume()}.

\begin{figure}[htbp]\centering
\begin{tikzpicture}[font=\small,
  nd/.style={draw,rounded corners=2pt,align=center,text width=34mm,
    minimum height=11mm,inner sep=3pt}]
\node[nd] (sc) {\texttt{Scene}};
\node[nd,below=20mm of sc] (vb) {\texttt{VBox}};
\node[nd,below=20mm of vb] (bt) {\texttt{Button}\\[1pt]\scriptsize цель события};
\draw[-Latex] ([xshift=-10mm]sc.south) -- ([xshift=-10mm]vb.north)
  node[midway,left,font=\scriptsize,align=right] {захват:\\фильтры};
\draw[-Latex] ([xshift=-10mm]vb.south) -- ([xshift=-10mm]bt.north)
  node[midway,left,font=\scriptsize,align=right] {захват:\\фильтры};
\draw[-Latex] ([xshift=10mm]bt.north) -- ([xshift=10mm]vb.south)
  node[midway,right,font=\scriptsize,align=left] {всплытие:\\обработчики};
\draw[-Latex] ([xshift=10mm]vb.north) -- ([xshift=10mm]sc.south)
  node[midway,right,font=\scriptsize,align=left] {всплытие:\\обработчики};
\end{tikzpicture}
\caption{Две фазы распространения события по дереву узлов}\label{fig:09-5}
\end{figure}

Фильтры, добавленные методом \texttt{addEventFilter}, срабатывают в фазе
захвата~--- сверху вниз; обработчики (\texttt{addEventHandler} и все методы вида
\texttt{setOnXxx})~--- в фазе всплытия, снизу вверх. Поэтому если фильтр стоит
на контейнере, обработчик~--- на кнопке, а ещё один~--- снова на контейнере,
порядок будет такой: фильтр контейнера, обработчик кнопки, обработчик
контейнера. Вызов \texttt{event.consume()} в фильтре лишает кнопку события
вовсе: так централизованно блокируют ввод на время загрузки данных. И ещё одно
отличие: \texttt{setOnAction} у узла один, новый вызов затирает предыдущий, а
\texttt{addEventHandler} можно вызывать сколько угодно раз.

\subsection{Свойства и привязки}

Обычное поле \texttt{String name} умеет только хранить значение.
\term{Свойство} (\eng{property}) умеет хранить значение и сообщать всем
заинтересованным, что оно изменилось. Аналогия: обычное поле~--- листок с
записанным курсом валюты; свойство~--- табло в обменнике: как только курс
поменялся, все, кто на него смотрит, узнают об этом мгновенно.

Типы свойств повторяют базовые типы данных: \texttt{StringProperty},
\texttt{IntegerProperty}, \texttt{DoubleProperty}, \texttt{BooleanProperty} и
\texttt{ObjectProperty<T>}; у каждого есть готовая реализация с приставкой
\texttt{Simple}. Настоящая сила свойств~--- в
привязках: метод \texttt{bind()} создаёт одностороннюю связь, при которой
значение приёмника всегда повторяет значение источника, а сам приёмник
становится доступен только для чтения. Что можно выразить одними привязками,
показано в листинге~\ref{lst:09-10}: ни одного обработчика событий здесь нет.

\begin{listing}[H]
\begin{minted}{java}
// Метка всегда показывает то же, что набрано в поле
label.textProperty().bind(field.textProperty());

// Привязка к вычисляемому выражению
counter.textProperty().bind(Bindings.concat(
        "Введено символов: ", field.textProperty().length()));

// Кнопка неактивна, пока поле пустое
submit.disableProperty().bind(field.textProperty().isEmpty());

// Аналог тернарного оператора в мире привязок
status.textProperty().bind(
        Bindings.when(field.textProperty().isEmpty())
                .then("Заполните поле")
                .otherwise("Готово к отправке"));
\end{minted}
\caption{Односторонние привязки}
\label{lst:09-10}
\end{listing}

\begin{oshibka}
\textbf{Частая ошибка: запись в привязанное свойство.} После вызова
\texttt{label.textProperty().bind(...)} обращение \texttt{label.setText(...)} не
поменяет текст, а выбросит \texttt{RuntimeException}: связанное значение
изменять нельзя. Чтобы снова управлять текстом вручную, связь надо разорвать
вызовом \texttt{unbind()}.
\end{oshibka}

Метод \texttt{bindBidirectional()} связывает два свойства как сообщающиеся
сосуды: изменили одно~--- изменилось второе, и наоборот; оба остаются доступны
для записи. Классический сценарий~--- поле формы, связанное с полем модели.
Плата за удобство~--- совпадение типов: односторонняя связь типы допускает
разные (через выражения), двусторонняя требует одинаковых.

\texttt{ObservableList<T>}~--- список, который сообщает о добавлении и удалении
элементов; на нём построены таблица, список и выпадающий список. Связали список
с таблицей один раз~--- и дальше работаете со списком как с обычной коллекцией,
а интерфейс обновляется сам. Это работает как табло вылетов в аэропорту:
диспетчер меняет расписание в одном месте, и все экраны обновляются сами~---
никто не бегает по залу с тряпкой и маркером. Создаются такие списки фабрикой
\texttt{FXCollections.observableArrayList(...)}.

Две производные коллекции превращают наблюдаемый список в основу живого поиска:
\texttt{FilteredList} показывает только подходящие элементы, а
\texttt{SortedList} даёт сортированное представление. Обе не копируют данные, а
служат «окном» в исходный список: изменился источник~--- тут же меняется и то,
что видно через окно. Вместе они дают поиск по мере набора текста плюс
сортировку по щелчку на заголовке столбца (листинг~\ref{lst:09-11}).

\begin{listing}[H]
\begin{minted}{java}
// Окно в исходный список; данные не копируются
FilteredList<Student> filtered =
        new FilteredList<>(service.getStudents(), student -> true);

searchField.textProperty().addListener((obs, was, now) -> {
    String query = now == null ? "" : now.trim().toLowerCase();
    filtered.setPredicate(student -> query.isEmpty()
            || student.getFullName().toLowerCase().contains(query)
            || student.getGroup().toLowerCase().contains(query));
});

// Сортировка по щелчку на заголовке столбца
SortedList<Student> sorted = new SortedList<>(filtered);
sorted.comparatorProperty().bind(table.comparatorProperty());
table.setItems(sorted);
\end{minted}
\caption{Живой поиск и сортировка в таблице}
\label{lst:09-11}
\end{listing}

Последняя привязка обязательна: без неё щелчок по заголовку переключает
компаратор самой таблицы, а \texttt{SortedList} о нём ничего не знает и
продолжает отдавать элементы в прежнем порядке~--- сортировка на экране просто
не работает.

\section{Шаблон MVC: модель и сервис}

\term{MVC} (\eng{Model-View-Controller})~--- архитектурный шаблон, разделяющий
приложение на три части с чётко разными обязанностями. Аналогия~--- ресторан.
Кухня знает, из чего готовятся блюда, и понятия не имеет, как выглядит зал: это
\term{модель}. Зал~--- столы, скатерти, меню~--- это \term{представление}, он
ничего не готовит. Официант ходит между залом и кухней: принял заказ, передал на
кухню, принёс готовое блюдо~--- это \term{контроллер}. Поменяли интерьер зала~---
кухню переделывать не нужно; сменили рецепт~--- официант работает по-прежнему.
Как эти роли выглядят в приложении JavaFX, показывает табл.~\ref{tab:09-5}.

\begin{table}[htbp]
\caption{Роли шаблона MVC в приложении JavaFX}
\label{tab:09-5}
\centering
\begin{tabular}{L{3.3cm}L{4.0cm}L{3.9cm}L{3.6cm}}
\toprule
Роль & Что это в JavaFX & Обязанности & Чего не должна знать \\
\midrule
Модель & классы со свойствами, наблюдаемые списки, сервисы & данные, проверки, вычисления, работа с базой & ничего о кнопках, полях и разметке \\
Представление & файл FXML и таблица стилей & как выглядит окно и где что стоит & ничего о бизнес-правилах \\
Контроллер & класс с полями \texttt{@FXML} и обработчиками & реагирует на действия, вызывает модель & как модель хранит данные \\
\bottomrule
\end{tabular}
\end{table}

Поток управления при нажатии кнопки показан на рис.~\ref{fig:09-6}. Самая
красивая часть~--- последний шаг: контроллеру не нужно перерисовывать таблицу
вручную, он изменил модель, а привязка данных сделала остальное.

\begin{figure}[htbp]\centering
\begin{tikzpicture}[font=\small,node distance=5mm,
  st/.style={draw,rounded corners=2pt,align=center,text width=88mm,
    inner sep=4pt,minimum height=9mm}]
\node[st] (a) {Пользователь нажал кнопку в окне (представление, FXML)};
\node[st,below=of a] (b) {Разметка вызвала метод контроллера, указанный в
  атрибуте \texttt{onAction}};
\node[st,below=of b] (c) {Контроллер считал значения полей и проверил
  корректность};
\node[st,below=of c] (d) {Контроллер вызвал сервис: модель создала объект и
  добавила его в наблюдаемый список};
\node[st,below=of d] (e) {Список сообщил подписчикам об изменении};
\node[st,below=of e] (f) {Таблица обновилась сама};
\draw[flowarrow] (a) -- (b);
\draw[flowarrow] (b) -- (c);
\draw[flowarrow] (c) -- (d);
\draw[flowarrow] (d) -- (e);
\draw[flowarrow] (e) -- (f);
\end{tikzpicture}
\caption{Поток управления по шаблону MVC}\label{fig:09-6}
\end{figure}

Четыре практических правила удерживают архитектуру от расползания. Первое:
модель не импортирует \texttt{javafx.scene}~--- появился
\texttt{import javafx.scene.control.Label} в классе \texttt{Student}, и
архитектура сломана; импорт \texttt{javafx.beans.property} допустим, потому что
свойства относятся к модулю \texttt{javafx.base}. Второе: контроллер тонкий~---
вычисления и проверки живут в сервисе. Третье: одно окно~--- один контроллер.
Четвёртое: данные между окнами передаются через контроллер, полученный из
загрузчика, а не через статические поля.

Теперь соберём всё изученное в одно работающее приложение: таблица студентов,
форма добавления, удаление выбранной записи, счётчик и таблица стилей. Начинаем
с модели (листинг~\ref{lst:09-12}); устойчивый шаблон здесь такой: на каждое
поле~--- три метода, получатель, установщик и метод, возвращающий само
свойство.

\begin{listing}[H]
\begin{minted}{java}
package com.example.manager.model;

import javafx.beans.property.*;

public class Student {

    private final StringProperty fullName =
            new SimpleStringProperty(this, "fullName", "");
    private final StringProperty group =
            new SimpleStringProperty(this, "group", "");
    private final DoubleProperty averageGrade =
            new SimpleDoubleProperty(this, "averageGrade", 0.0);

    public Student(String fullName, String group, double grade) {
        this.fullName.set(fullName);
        this.group.set(group);
        this.averageGrade.set(grade);
    }

    public String getFullName() { return fullName.get(); }
    public void setFullName(String v) { fullName.set(v); }
    public StringProperty fullNameProperty() { return fullName; }

    public String getGroup() { return group.get(); }
    public void setGroup(String v) { group.set(v); }
    public StringProperty groupProperty() { return group; }

    public double getAverageGrade() { return averageGrade.get(); }
    public void setAverageGrade(double v) { averageGrade.set(v); }
    public DoubleProperty averageGradeProperty() {
        return averageGrade;
    }

    @Override
    public String toString() {
        return fullName.get() + " (" + group.get() + ")";
    }
}
\end{minted}
\caption{Модель: класс Student}
\label{lst:09-12}
\end{listing}

Обратите внимание: поля-свойства объявлены как \texttt{final}, хотя значения
меняются. Противоречия здесь нет~--- неизменна ссылка на объект-свойство, а
значение внутри него меняется методом \texttt{set}; именно поэтому все
подписчики и привязки остаются в силе на всё время жизни объекта.

Модель в нашем примере~--- не только класс \texttt{Student}, но и сервис вокруг
него (листинг~\ref{lst:09-13}). Он хранит единственный наблюдаемый список и
предоставляет операции над ним, чтобы ни контроллер, ни разметка не работали со
списком напрямую. Обратите внимание: здесь нет ни одного импорта
\texttt{javafx.scene}~--- это и есть признак правильно отделённой модели.

\begin{listing}[H]
\begin{minted}{java}
package com.example.manager.service;

import com.example.manager.model.Student;
import javafx.collections.FXCollections;
import javafx.collections.ObservableList;

public class StudentService {

    private final ObservableList<Student> students =
            FXCollections.observableArrayList();

    public StudentService() {
        students.addAll(
                new Student("Петров Иван", "ПИ24-1", 4.6),
                new Student("Смирнова Анна", "ПИ24-1", 4.9),
                new Student("Кузнецов Дмитрий", "ТРПО24-1", 3.8));
    }

    /** Список для привязки к таблице. */
    public ObservableList<Student> getStudents() {
        return students;
    }

    public void add(Student student) { students.add(student); }

    public void remove(Student student) { students.remove(student); }

    public static boolean isValidGrade(double grade) {
        return grade >= 2.0 && grade <= 5.0;
    }
}
\end{minted}
\caption{Модель: класс StudentService}
\label{lst:09-13}
\end{listing}

\section{FXML, контроллер и оформление}

Всё, что мы писали до сих пор, строило интерфейс кодом. Для окна из пяти
элементов это нормально; для окна из пятидесяти метод \texttt{start()}
превращается в километровую простыню, в которой невозможно найти нужную кнопку.
Аналогия: прораб может объяснять бригаде каждый кирпич устно, но проще один раз
начертить план~--- тогда чертёж отдельно (его правит архитектор), а работа
отдельно (её делает бригада).

\term{FXML}~--- XML-формат описания графа сцены. Он даёт разделение
ответственности, читаемость, возможность рисовать интерфейс мышью и правку
разметки без перекомпиляции. Разметка нашего окна приведена в
листинге~\ref{lst:09-14}: сравните её с деревом на рис.~\ref{fig:09-3}~---
вложенность элементов XML буквально повторяет вложенность узлов. Строки формы
однотипны, поэтому показана одна из трёх: поля «Группа» и «Средний балл»
устроены точно так же и отличаются только именем и номером строки.

\begin{listing}[H]
\begin{minted}{xml}
<?xml version="1.0" encoding="UTF-8"?>

<?import javafx.geometry.Insets?>
<?import javafx.scene.control.*?>
<?import javafx.scene.layout.*?>

<BorderPane xmlns="http://javafx.com/javafx/21"
    xmlns:fx="http://javafx.com/fxml/1"
    fx:controller="com.example.manager.controller.StudentController"
    stylesheets="@styles.css" prefWidth="780" prefHeight="540">

  <top>
    <Label text="Менеджер студентов" styleClass="title-label"/>
  </top>

  <center>
    <TableView fx:id="studentTable">
      <columns>
        <TableColumn fx:id="fullNameColumn" text="ФИО" prefWidth="330"/>
        <TableColumn fx:id="groupColumn" text="Группа" prefWidth="150"/>
        <TableColumn fx:id="gradeColumn" text="Балл" prefWidth="150"/>
      </columns>
    </TableView>
  </center>

  <bottom>
    <VBox spacing="12">
      <padding>
        <Insets top="14" right="18" bottom="18" left="18"/>
      </padding>

      <GridPane hgap="10" vgap="10">
        <!-- Статические свойства панели: GridPane.rowIndex -->
        <Label text="ФИО:" styleClass="field-label"
               GridPane.columnIndex="0" GridPane.rowIndex="0"/>
        <TextField fx:id="fullNameField" prefColumnCount="24"
               GridPane.columnIndex="1" GridPane.rowIndex="0"/>
      </GridPane>

      <HBox spacing="10">
        <!-- onAction ссылается на метод контроллера; # обязателен -->
        <Button fx:id="addButton" text="Добавить"
                onAction="#handleAdd" styleClass="primary-button"/>
        <Button fx:id="deleteButton" text="Удалить выбранного"
                onAction="#handleDelete" styleClass="danger-button"/>
        <Button text="Очистить форму" onAction="#handleClear"/>
      </HBox>

      <Label fx:id="statusLabel" styleClass="status-label"/>
    </VBox>
  </bottom>
</BorderPane>
\end{minted}
\caption{Представление: файл students-view.fxml}
\label{lst:09-14}
\end{listing}

Каждой конструкции разметки механически соответствует строка кода на Java:
\texttt{<VBox spacing="12">} равносильно \texttt{vbox.setSpacing(12)}, атрибут
\texttt{GridPane.rowIndex="1"}~--- вызову
\texttt{GridPane.setRowIndex(node, 1)}. Три атрибута пространства имён
\texttt{fx} нужно знать наизусть. \texttt{fx:controller}~--- полное имя
класса-контроллера, указывается только на корневом элементе.
\texttt{fx:id}~--- идентификатор узла: по нему загрузчик находит одноимённое
поле контроллера и подставляет туда созданный объект.
\texttt{fx:include}~--- вставка другого файла разметки.

\subsection{Контроллер и загрузчик}

Класс-контроллер связывает разметку с поведением: он объявляет поля с аннотацией
\texttt{@FXML}, имена которых совпадают с \texttt{fx:id}, и методы, на которые
ссылается \texttt{onAction}. Начнём с настройки (листинг~\ref{lst:09-15}): в ней
собрано почти всё, о чём шла речь в главе,~--- фабрики значений, связь таблицы
со списком и привязки, благодаря которым кнопки и строка состояния управляют
собой сами.

\begin{listing}[H]
\begin{minted}{java}
package com.example.manager.controller;

// импорты: Student, StudentService, Bindings, FXML,
// javafx.scene.control.*, PropertyValueFactory

public class StudentController {

    // Имена полей совпадают с fx:id в разметке
    @FXML private TableView<Student> studentTable;
    @FXML private TableColumn<Student, String> fullNameColumn;
    @FXML private TableColumn<Student, String> groupColumn;
    @FXML private TableColumn<Student, Double> gradeColumn;
    @FXML private TextField fullNameField, groupField, gradeField;
    @FXML private Button addButton, deleteButton;
    @FXML private Label statusLabel;

    private final StudentService service = new StudentService();

    /** Вызывается ПОСЛЕ заполнения всех полей @FXML. */
    @FXML
    private void initialize() {
        // 1. Откуда столбцы берут значения
        fullNameColumn.setCellValueFactory(
                new PropertyValueFactory<>("fullName"));
        groupColumn.setCellValueFactory(
                new PropertyValueFactory<>("group"));
        gradeColumn.setCellValueFactory(
                d -> d.getValue().averageGradeProperty().asObject());

        // 2. Связь таблицы со списком модели
        studentTable.setItems(service.getStudents());
        studentTable.setPlaceholder(new Label("Список пуст"));

        // 3. Кнопки включаются и выключаются сами
        deleteButton.disableProperty().bind(studentTable
                .getSelectionModel().selectedItemProperty().isNull());
        addButton.disableProperty().bind(
                fullNameField.textProperty().isEmpty()
                        .or(groupField.textProperty().isEmpty())
                        .or(gradeField.textProperty().isEmpty()));

        // 4. Строка состояния показывает количество записей
        statusLabel.textProperty().bind(Bindings.concat(
                "Студентов в списке: ",
                Bindings.size(service.getStudents())));
    }
}
\end{minted}
\caption{Контроллер: настройка в методе initialize}
\label{lst:09-15}
\end{listing}

Порядок работы загрузчика \texttt{FXMLLoader} полезно знать целиком: он читает
XML и создаёт объекты узлов; создаёт экземпляр контроллера из
\texttt{fx:controller} (тому нужен конструктор без аргументов); записывает в
поля с \texttt{@FXML} узлы с соответствующими \texttt{fx:id}; привязывает
обработчики; вызывает \texttt{initialize()}, если он есть; возвращает корневой
узел.

\begin{oshibka}
\textbf{Частая ошибка: настройка в конструкторе контроллера.} В конструкторе
поля \texttt{@FXML} ещё равны \texttt{null}: загрузчик заполняет их позже.
Обращение к ним оттуда даёт \texttt{NullPointerException}. Вся настройка
интерфейса идёт в методе \texttt{initialize()}.
\end{oshibka}

Осталось написать обработчики кнопок (листинг~\ref{lst:09-16}). Разбор оценки
вынесен в отдельный метод, чтобы не дублировать его при добавлении и изменении
записи, а об ошибках пользователю сообщает диалог, а не консоль. Не показаны два
коротких метода: \texttt{handleClear} снимает выделение и очищает поля формы, а
\texttt{showError} показывает \texttt{Alert} типа \texttt{ERROR}.

\begin{listing}[H]
\begin{minted}{java}
    @FXML
    private void handleAdd() {
        Double grade = parseGrade();
        if (grade == null) {
            return;
        }
        service.add(new Student(fullNameField.getText().trim(),
                groupField.getText().trim(), grade));
        handleClear();
    }

    @FXML
    private void handleDelete() {
        Student selected =
                studentTable.getSelectionModel().getSelectedItem();
        if (selected == null) {
            return;   // кнопка и так отключена, но защита не мешает
        }
        Alert confirm = new Alert(Alert.AlertType.CONFIRMATION);
        confirm.setHeaderText(null);
        confirm.setContentText("Удалить запись: " + selected + "?");
        confirm.showAndWait()
                .filter(button -> button == ButtonType.OK)
                .ifPresent(button -> service.remove(selected));
    }

    /** Разбор поля оценки; null, если значение некорректно. */
    private Double parseGrade() {
        double grade;
        try {
            grade = Double.parseDouble(
                    gradeField.getText().trim().replace(',', '.'));
        } catch (NumberFormatException e) {
            showError("Балл должен быть числом, например 4.5");
            return null;
        }
        if (!StudentService.isValidGrade(grade)) {
            showError("Балл должен быть от 2.0 до 5.0");
            return null;
        }
        return grade;
    }
\end{minted}
\caption{Контроллер: обработчики кнопок}
\label{lst:09-16}
\end{listing}

Загружают разметку двумя способами. Короткий~--- это статический вызов
\texttt{FXMLLoader.load(url)}, но он возвращает только корневой узел, и
контроллер у него не получить. Поэтому по умолчанию создают объект загрузчика,
как в точке входа (листинг~\ref{lst:09-17}), а контроллер добирают методом
\texttt{getController()}. Обёртка \texttt{Objects.requireNonNull} здесь не
формальность: без неё вместо внятного сообщения о пропавшем файле вы получите
загадочное \texttt{Location is required}.

\begin{listing}[H]
\begin{minted}{java}
package com.example.manager;

// импорты: Application, FXMLLoader, Parent, Scene, Stage,
// java.net.URL, java.util.Objects

public class MainApp extends Application {

    @Override
    public void start(Stage stage) throws Exception {
        // Ресурс ищется рядом с этим классом в resources
        URL view = MainApp.class.getResource("students-view.fxml");
        Objects.requireNonNull(view, "Не найден students-view.fxml");

        Parent root = new FXMLLoader(view).load();

        stage.setTitle("Менеджер студентов");
        stage.setScene(new Scene(root, 780, 540));
        stage.show();
    }

    public static void main(String[] args) {
        launch(args);
    }
}
\end{minted}
\caption{Точка входа: класс MainApp}
\label{lst:09-17}
\end{listing}

\subsection{Scene Builder и оформление через CSS}

\term{JavaFX Scene Builder}~--- бесплатный визуальный редактор интерфейсов от
компании Gluon: вы перетаскиваете компоненты мышью, а он сохраняет результат в
файл FXML, и наоборот~--- любая корректная разметка открывается в нём как
готовое окно. Ключевая мысль: Java-кода он не генерирует и средой разработки не
является, он редактирует ровно один тип файлов; \texttt{fx:id} и имя
метода-обработчика задают в его вкладке \eng{Code}. Помогает редактор, когда
интерфейс большой и статичный; мешает~--- когда набор элементов формируется
динамически и когда файл правят несколько человек, потому что он переставляет
атрибуты и слияние версий превращается в мучение. Отсюда компромисс, принятый в
индустрии: статический каркас окна~--- в FXML, динамические части~--- кодом.

Задавать цвет каждой кнопке вызовом \texttt{setStyle()}~--- то же самое, что
пришивать пуговицы к каждой рубашке отдельно, вместо того чтобы один раз описать
фасон. Таблица стилей даёт единое место, где хранится весь внешний вид: поменяли
одну строку~--- изменилось всё приложение. По умолчанию JavaFX использует
встроенную тему \eng{Modena}; ваша таблица стилей её не заменяет, а
накладывается сверху.

Способов применить стиль три: инлайн-стиль \texttt{setStyle("...")}~--- быстро,
но не переиспользуется; класс стиля плюс внешняя таблица
(\texttt{getStyleClass().add("primary-button")})~--- основной рабочий способ;
идентификатор \texttt{setId("submitButton")}~--- для уникального элемента. Сама
таблица подключается к сцене через список \texttt{getStylesheets()}, в который
добавляют адрес файла; вызов \texttt{toExternalForm()} обязателен, потому что
список хранит строки-адреса, а \texttt{getResource()} возвращает объект
\texttt{URL}. Можно подключить таблицу и из разметки атрибутом
\texttt{stylesheets="@styles.css"}, где знак собачки означает «путь относительно
этого файла»,~--- так сделано в листинге~\ref{lst:09-14}.

Таблица стилей нашего приложения приведена в листинге~\ref{lst:09-18}. Разберите
её по типам селекторов: точка~--- класс стиля, знак номера~--- идентификатор, а
имя после двоеточия~--- псевдокласс состояния.

\begin{listing}[H]
\begin{minted}{text}
/* Корневой узел сцены всегда имеет класс .root */
.root {
    -fx-font-family: "Segoe UI", "Arial", sans-serif;
    -fx-font-size: 13px;
    -fx-background-color: #f4f6f8;
}

.title-label  { -fx-font-size: 22px; -fx-font-weight: bold; }
.field-label  { -fx-text-fill: #55606d; }
.status-label { -fx-text-fill: #6b7684; -fx-font-size: 12px; }

.text-field { -fx-border-color: #d3dae3; -fx-padding: 6 10 6 10; }
.text-field:focused { -fx-border-color: #2d7ff9; }

.button { -fx-background-radius: 5; -fx-cursor: hand; }

.primary-button {
    -fx-background-color: #2d7ff9;
    -fx-text-fill: white;
    -fx-font-weight: bold;
}
.primary-button:hover { -fx-background-color: derive(#2d7ff9, -12%); }
.primary-button:disabled { -fx-opacity: 0.45; }

.danger-button { -fx-background-color: #e5484d; -fx-text-fill: white; }

.table-view .column-header { -fx-background-color: #eaeef3; }
\end{minted}
\caption{Таблица стилей styles.css}
\label{lst:09-18}
\end{listing}

Имя селектора по типу компонента выводится из имени класса: \texttt{Button}~---
\texttt{.button}, \texttt{TextField}~--- \texttt{.text-field}. Основные
псевдоклассы: \texttt{:hover}, \texttt{:pressed}, \texttt{:focused},
\texttt{:disabled}, \texttt{:selected}, а также \texttt{:odd} и \texttt{:even}
для строк списков. Синтаксис почти совпадает с веб-вариантом CSS, но набор
свойств свой, и все они начинаются с приставки \texttt{-fx-}. Ключевые
расхождения собраны в табл.~\ref{tab:09-6}; особое внимание обратите на две
последние строки~--- именно их ищут в первую очередь и не находят.

\begin{table}[htbp]
\caption{Отличия CSS в JavaFX от CSS в вебе}
\label{tab:09-6}
\centering
\begin{tabular}{L{4.3cm}L{5.5cm}L{4.9cm}}
\toprule
В вебе & В JavaFX & Комментарий \\
\midrule
\texttt{color} & \texttt{-fx-text-fill} & цвет текста \\
\texttt{background-color} & \texttt{-fx-background-color} & допускает несколько слоёв \\
\texttt{border-radius} & \texttt{-fx-background-radius} и \texttt{-fx-border-radius} & фон и рамка скругляются отдельно \\
\texttt{box-shadow} & \texttt{-fx-effect: dropshadow(...)} & тень~--- это эффект \\
\texttt{margin} & отсутствует & внешний отступ задаёт панель \\
\texttt{display}, \texttt{flex} & отсутствуют & раскладка~--- задача панелей \\
\bottomrule
\end{tabular}
\end{table}

Запомните и приоритет: инлайн-стиль перебивает вашу таблицу стилей, а та~---
встроенную тему, поэтому не смешивайте два подхода для одного свойства.
Запускается готовое приложение командой \texttt{mvn clean javafx:run}, а из
среды разработки~--- классом \texttt{Launcher} из листинга~\ref{lst:09-2}.
Развивать пример можно в любую сторону: добавить поиск через
\texttt{FilteredList} или подключить базу данных из главы~\ref{ch:06}. Модель
при этом почти не изменится~--- в этом и состоит выгода разделения слоёв.

\praktikum{}

Понадобятся JDK~21 или новее, Maven~3.8 или новее и любая среда разработки.
Задания выполняются в одном проекте и строго по порядку: каждое следующее
опирается на файлы, созданные в предыдущем.

\subsection*{Задание 1. Проект, первое окно и жизненный цикл}

Создайте каталог \texttt{student-manager} со структурой из рис.~\ref{fig:09-1};
сборочный файл возьмите из листинга~\ref{lst:09-1}, первое окно~--- из
листинга~\ref{lst:09-3}. Запустите проект командой
\texttt{mvn clean compile javafx:run}.

Добавьте в \texttt{HelloApp} методы \texttt{init()} и \texttt{stop()},
печатающие имя текущего потока (\texttt{Thread.currentThread().getName()}), и
такую же печать в начало \texttt{start()}. Добавьте кнопку «Выйти» с
обработчиком \texttt{Platform.exit()}. Закройте окно кнопкой и системным
крестиком, сравните вывод консоли, затем замените \texttt{Platform.exit()} на
\texttt{System.exit(0)} и проверьте, вызывается ли \texttt{stop()}. Наконец,
запустите \texttt{HelloApp} прямо из среды разработки: скорее всего вы получите
сообщение о недостающих компонентах JavaFX; создайте класс \texttt{Launcher} по
листингу~\ref{lst:09-2} и запустите его.

Ответьте письменно: в каких потоках выполнились \texttt{init()} и
\texttt{start()}; вызывается ли \texttt{stop()} при каждом из трёх способов
завершения; почему \texttt{Launcher} решает проблему, запуская тот же класс.

\subsection*{Задание 2. Элементы управления, обработчики и привязки}

Создайте класс \texttt{GreetingApp}: подпись, поле ввода и кнопки
«Поздороваться» (\texttt{setDefaultButton(true)}) и «Очистить»
(\texttt{setCancelButton(true)}). Обработчик первой кнопки читает фамилию из
поля и выводит приветствие в метку результата, а если поле пустое~--- сообщает
об этом; тот же обработчик повесьте на \texttt{setOnAction} самого поля.
Разложите элементы по \texttt{HBox} внутри \texttt{VBox} по образцу
листинга~\ref{lst:09-5}.

Вторая часть~--- то же самое без единого обработчика. Добавьте три метки и
настройте их привязками по образцу листинга~\ref{lst:09-10}: первая повторяет
содержимое поля, вторая показывает число введённых символов, третья сообщает
«Заполните поле» или «Можно здороваться». Кнопку свяжите с пустотой поля через
\texttt{disableProperty()}. Затем попробуйте задать текст первой метки вызовом
\texttt{setText}, запишите текст исключения и уберите эту строку.

Ответьте письменно: чем \texttt{setDefaultButton(true)} отличается от
\texttt{setOnAction}; какое исключение дала запись в привязанное свойство и
почему.

\subsection*{Задание 3. Компоновка, модель и таблица}

Создайте класс \texttt{FormApp} с формой регистрации: сетка «подпись~--- поле»
из листинга~\ref{lst:09-6} с полями «Логин», «Пароль» (\texttt{PasswordField}) и
«Группа» (\texttt{ComboBox}), под ней~--- строка кнопок, ниже~--- строка
состояния; всё это в \texttt{VBox}, а его~--- в центр \texttt{BorderPane} с
заголовком сверху. Растяните окно мышью: поля должны растягиваться, подписи~---
оставаться на месте. Проверьте два изменения: включите
\texttt{setGridLinesVisible(true)}; удалите обе строки с
\texttt{ColumnConstraints} и снова растяните окно.

Вторая часть~--- данные. Наберите классы \texttt{Student} и
\texttt{StudentService} из листингов~\ref{lst:09-12} и~\ref{lst:09-13}, затем
класс \texttt{TableApp} с таблицей по образцу листингов~\ref{lst:09-8}
и~\ref{lst:09-9}. Добавьте кнопку, добавляющую студента в сервис, и кнопку
удаления выбранной строки, связав вторую привязкой с пустым выделением; строку
состояния сделайте привязкой к \texttt{Bindings.size(...)}. Затем внесите
опечатку в имя свойства внутри \texttt{PropertyValueFactory}, пересоберите
проект, посмотрите на результат и верните имя обратно.

Ответьте письменно: чем внутренний отступ отличается от внешнего; зачем нужен
\texttt{Priority.ALWAYS}; скомпилировался ли проект с опечаткой и почему
компилятор промолчал; кто перерисовал таблицу после добавления студента.

\subsection*{Задание 4. Разметка FXML, контроллер и таблица стилей}

Вынесите интерфейс в разметку. Создайте файл \texttt{students-view.fxml} по
листингу~\ref{lst:09-14}, дописав третью строку формы с полем
\texttt{gradeField}; контроллер \texttt{StudentController} по
листингам~\ref{lst:09-15} и~\ref{lst:09-16}, добавив метод
\texttt{handleClear}; точку входа \texttt{MainApp} по листингу~\ref{lst:09-17} и
таблицу стилей \texttt{styles.css} по листингу~\ref{lst:09-18}.

Проведите три эксперимента. Первый: переименуйте в разметке
\texttt{fx:id="groupField"} в \texttt{fx:id="groupFld"}, оставив поле
контроллера прежним, и запишите текст исключения. Второй: перенесите настройку
столбцов из \texttt{initialize()} в конструктор контроллера и объясните
результат. Третий: добавьте в конец \texttt{initialize()} вызов
\texttt{addButton.setStyle}, задающий зелёный фон, и посмотрите, какой цвет
победил. Заодно проверьте псевдоклассы: наведите курсор на синюю кнопку,
щёлкните по полю ввода, очистите форму, выделите строку таблицы.

Ответьте письменно: какое исключение дало несовпадение \texttt{fx:id} и имени
поля; почему настройка не работает из конструктора; как расставить по приоритету
инлайн-стиль, вашу таблицу стилей и встроенную тему.

\subsection*{Задание 5. Итоговое приложение: изменение, поиск, подтверждения}

Добавьте в разметку кнопку «Сохранить изменения» с обработчиком
\texttt{handleUpdate}, а в контроллер~--- одноимённый метод: он берёт выбранного
студента, разбирает поля формы методом \texttt{parseGrade()} и обновляет
свойства объекта. Добавьте в \texttt{initialize()} слушателя выбранной строки,
заполняющего форму её значениями, и привяжите доступность новой кнопки к наличию
выделения.

Добавьте живой поиск: поле \texttt{searchField} в верхнюю зону разметки и блок
из листинга~\ref{lst:09-11} вместо вызова \texttt{setItems}; строку состояния
свяжите сразу с двумя числами~--- сколько студентов всего и сколько показано.
Добавьте в \texttt{MainApp} обработчик \texttt{stage.setOnCloseRequest}, который
показывает диалог и вызывает \texttt{event.consume()}, если пользователь
передумал. Проверьте: ввод слова вместо числа даёт диалог с ошибкой, а не
падение; оценка 7,5 отвергается проверкой; поиск фильтрует таблицу по мере
набора.

Ответьте письменно: почему таблица обновилась после изменения свойств студента,
хотя вы не вызывали ни \texttt{refresh()}, ни \texttt{setItems()}; почему при
удалении отфильтрованной строки удаляется правильный студент; что делает
\texttt{event.consume()} в обработчике закрытия окна; есть ли в классах
\texttt{Student} и \texttt{StudentService} импорт \texttt{javafx.scene} и почему
его отсутствие~--- хороший признак.

\subsection*{Результаты занятия}

По итогам занятия у вас должны быть готовы: проект \texttt{student-manager} с
зависимостями JavaFX и классом \texttt{Launcher}; приложения \texttt{HelloApp},
\texttt{GreetingApp}, \texttt{FormApp} и \texttt{TableApp}; модель
\texttt{Student} и сервис \texttt{StudentService}; разметка
\texttt{students-view.fxml}, контроллер \texttt{StudentController} и таблица
стилей \texttt{styles.css}; итоговое приложение с добавлением, изменением,
удалением с подтверждением и поиском; снимки экрана итогового окна; письменные
ответы на вопросы заданий.

\kontrolvoprosy

\begin{enumerate}
\item Почему начиная с JDK~11 JavaFX не входит в состав JDK и какие два модуля
      достаточно объявить для большинства приложений?
\item Что означает ошибка о недостающих компонентах среды выполнения JavaFX и
      почему её обходит класс-запускалка?
\item Опишите жизненный цикл приложения. В каких потоках выполняются
      \texttt{init()} и \texttt{start()}, почему в \texttt{init()} нельзя
      создавать окно и чем \texttt{Platform.exit()} отличается от
      \texttt{System.exit(0)}?
\item Что такое JavaFX Application Thread и почему нельзя менять интерфейс из
      другого потока? Как вернуть в интерфейс результат фоновой работы?
\item Чем \texttt{Stage} отличается от \texttt{Scene}? Может ли одна и та же
      сцена стоять в двух окнах? Чем \texttt{show()} отличается от
      \texttt{showAndWait()}?
\item Что такое граф сцены? Чем различаются \texttt{Node}, \texttt{Parent},
      \texttt{Region} и \texttt{Control}?
\item Что произойдёт, если один и тот же узел добавить сначала в одну панель, а
      затем в другую? Как в JavaFX определяется порядок наложения узлов?
\item В чём разница между внутренним и внешним отступом? Какую панель вы
      выберете для формы, а какую~--- для каркаса окна?
\item Из каких трёх частей состоит работающий \texttt{TableView}? За что
      отвечает фабрика значений, а за что фабрика ячеек? Какую ошибку не
      поймает компилятор при использовании \texttt{PropertyValueFactory}?
\item В каком порядке срабатывают фильтры и обработчики события и что делает
      \texttt{consume()}? Чем \texttt{setOnAction} отличается от
      \texttt{addEventHandler}?
\item Чем \texttt{bind()} отличается от \texttt{bindBidirectional()}? Что
      произойдёт при попытке задать текст метки с привязанным свойством?
\item Что делают \texttt{fx:controller}, \texttt{fx:id} и \texttt{onAction} в
      разметке? Почему настройку компонентов пишут в \texttt{initialize()}, а не
      в конструкторе контроллера?
\item Назовите три отличия CSS в JavaFX от CSS в вебе. В каком порядке
      применяются инлайн-стиль, ваша таблица стилей и встроенная тема?
\item Как распределяются обязанности между моделью, представлением и
      контроллером? Какой импорт в классе модели сигнализирует о нарушении
      архитектуры?
\end{enumerate}

\testzadaniya

\begin{enumerate}

\item Начиная с какой версии JDK библиотека JavaFX не поставляется вместе с JDK?
  \begin{enumerate}
    \item[а)] JDK~8
    \item[б)] JDK~9
    \item[в)] JDK~11
    \item[г)] JDK~17
  \end{enumerate}

\item Какие два модуля достаточно объявить для приложения с обычными элементами
      управления и разметкой FXML?
  \begin{enumerate}
    \item[а)] \texttt{javafx-base} и \texttt{javafx-media}
    \item[б)] \texttt{javafx-graphics} и \texttt{javafx-web}
    \item[в)] все семь модулей, иначе приложение не запустится
    \item[г)] \texttt{javafx-controls} и \texttt{javafx-fxml}
  \end{enumerate}

\item В каком порядке выполняются шаги жизненного цикла приложения?
  \begin{enumerate}
    \item[а)] \texttt{init} $\to$ \texttt{launch} $\to$ \texttt{start} $\to$ \texttt{stop}
    \item[б)] \texttt{launch} $\to$ \texttt{start} $\to$ \texttt{init} $\to$ \texttt{stop}
    \item[в)] \texttt{start} $\to$ \texttt{init} $\to$ \texttt{stop} $\to$ \texttt{launch}
    \item[г)] \texttt{launch} $\to$ \texttt{init} $\to$ \texttt{start} $\to$ \texttt{stop}
  \end{enumerate}

\item Чем \texttt{System.exit(0)} отличается от \texttt{Platform.exit()}?
  \begin{enumerate}
    \item[а)] ничем: оба корректно завершают среду выполнения
    \item[б)] \texttt{Platform.exit()} закрывает только текущее окно
    \item[в)] \texttt{System.exit(0)} обрывает работу JVM немедленно, поэтому
              \texttt{stop()} не выполняется
    \item[г)] \texttt{System.exit(0)} вызывает \texttt{stop()}, а
              \texttt{Platform.exit()}~--- нет
  \end{enumerate}

\item Фоновый поток напрямую задаёт текст метки вызовом \texttt{setText}. Что
      произойдёт?
  \begin{enumerate}
    \item[а)] текст обновится: платформа сама переключится в нужный поток
    \item[б)] будет выброшено \texttt{IllegalStateException}
    \item[в)] ничего не произойдёт, вызов будет молча проигнорирован
    \item[г)] приложение зависнет до завершения фонового потока
  \end{enumerate}

\item Каким должен быть корневой узел, передаваемый в конструктор
      \texttt{Scene}?
  \begin{enumerate}
    \item[а)] любым наследником \texttt{Node}
    \item[б)] обязательно объектом класса \texttt{Group}
    \item[в)] наследником \texttt{Parent}, способным иметь потомков
    \item[г)] обязательно объектом класса \texttt{AnchorPane}
  \end{enumerate}

\item Кнопку добавили сначала в \texttt{VBox}, затем её же в \texttt{HBox}. Что
      окажется на экране?
  \begin{enumerate}
    \item[а)] кнопка появится сразу в обеих панелях
    \item[б)] кнопка останется в \texttt{VBox}, второе добавление
              проигнорируется
    \item[в)] будет выброшено исключение о дублировании потомков
    \item[г)] кнопка окажется только в \texttt{HBox}~--- из \texttt{VBox} она
              исчезнет
  \end{enumerate}

\item Чем внутренний отступ отличается от внешнего?
  \begin{enumerate}
    \item[а)] внутренний задаётся в пикселях, внешний~--- в процентах
    \item[б)] внутренний~--- расстояние от краёв панели до потомков,
              внешний~--- место вокруг одного потомка, задаваемое статическим
              методом панели
    \item[в)] внутренний применяется к элементам управления, внешний~--- к
              панелям
    \item[г)] это синонимы: оба метода делают одно и то же
  \end{enumerate}

\item Кнопка лежит в \texttt{VBox}. На \texttt{VBox} добавлен фильтр, на кнопке
      задан \texttt{setOnAction}, на \texttt{VBox}~--- ещё и обработчик. В каком
      порядке они сработают?
  \begin{enumerate}
    \item[а)] обработчик кнопки, фильтр, обработчик \texttt{VBox}
    \item[б)] фильтр, обработчик \texttt{VBox}, обработчик кнопки
    \item[в)] фильтр, обработчик кнопки, обработчик \texttt{VBox}
    \item[г)] все три сработают одновременно, порядок не определён
  \end{enumerate}

\item Текст метки привязан к тексту поля вызовом \texttt{bind}. Что произойдёт
      при попытке задать текст метки вызовом \texttt{setText}?
  \begin{enumerate}
    \item[а)] метка покажет новый текст, а привязка автоматически разорвётся
    \item[б)] метка покажет новый текст, а затем вернётся к тексту поля
    \item[в)] вызов \texttt{setText} выбросит \texttt{RuntimeException}
    \item[г)] код не скомпилируется
  \end{enumerate}

\item Почему таблица обновляется сама, когда в модель добавили новую запись?
  \begin{enumerate}
    \item[а)] потому что таблица раз в секунду перечитывает источник данных
    \item[б)] потому что таблица связана с наблюдаемым списком, который сообщает
              подписчикам об изменениях
    \item[в)] потому что контроллер обязан вызывать \texttt{refresh()}
    \item[г)] потому что \texttt{PropertyValueFactory} следит за изменениями
              через рефлексию
  \end{enumerate}

\item Почему настройку интерфейса пишут в \texttt{initialize()}, а не в
      конструкторе контроллера?
  \begin{enumerate}
    \item[а)] потому что конструктор выполняется в другом потоке
    \item[б)] потому что в конструкторе поля \texttt{@FXML} ещё равны
              \texttt{null}: загрузчик заполняет их позже
    \item[в)] потому что конструктор вызывается один раз, а
              \texttt{initialize()}~--- при каждом показе окна
    \item[г)] потому что в конструкторе запрещено обращаться к классам модели
  \end{enumerate}

\item Какого свойства в CSS для JavaFX нет вовсе?
  \begin{enumerate}
    \item[а)] внешнего отступа: его задаёт панель компоновки
    \item[б)] \texttt{-fx-padding}
    \item[в)] \texttt{-fx-opacity}
    \item[г)] \texttt{-fx-background-color}
  \end{enumerate}

\end{enumerate}
